\documentclass[11pt,a4paper]{article}

\usepackage[top=0.85in, bottom=0.85in, left=0.8in, right=0.8in]{geometry}
\usepackage{mathptmx}
\usepackage{xcolor}
\usepackage{listings}
\usepackage{fancyhdr}
\usepackage{hyperref}

\setlength{\headheight}{14.5pt}
\pagestyle{fancy}
\fancyhf{}
\lhead{\textbf{Data Structures Lab}}
\rhead{Linked Queue}
\cfoot{Page \thepage}
\renewcommand{\headrulewidth}{0.6pt}

\hypersetup{
    colorlinks=true,
    linkcolor=blue,
    urlcolor=blue,
    pdftitle={Linked Queue Program - Student Records},
    pdfauthor={Data Structures Lab}
}

% Code listing styling
\definecolor{codebg}{rgb}{0.97,0.98,0.99}
\definecolor{codeframe}{rgb}{0.82,0.85,0.90}
\definecolor{keywordcolor}{rgb}{0.0,0.25,0.7}
\definecolor{commentcolor}{rgb}{0.38,0.45,0.52}
\definecolor{stringcolor}{rgb}{0.12,0.52,0.22}

\lstdefinestyle{cstyle}{
    language=C,
    basicstyle=\fontsize{9}{11.5}\selectfont\ttfamily,
    keywordstyle=\color{keywordcolor}\bfseries,
    commentstyle=\color{commentcolor}\itshape,
    stringstyle=\color{stringcolor},
    numbers=left,
    numberstyle=\tiny\color{gray},
    stepnumber=1,
    numbersep=8pt,
    backgroundcolor=\color{codebg},
    frame=single,
    rulecolor=\color{codeframe},
    tabsize=4,
    showspaces=false,
    showstringspaces=false,
    breaklines=true,
    breakatwhitespace=false,
    xleftmargin=12pt,
    framexleftmargin=12pt,
    captionpos=b
}

\lstdefinestyle{terminalstyle}{
    basicstyle=\fontsize{8.5}{11}\selectfont\ttfamily,
    backgroundcolor=\color{black!4},
    frame=single,
    rulecolor=\color{black!25},
    breaklines=true,
    breakatwhitespace=false,
    numbers=none,
    xleftmargin=8pt,
    framexleftmargin=8pt,
    tabsize=4
}

\begin{document}

\begin{center}
    {\Large \textbf{DATA STRUCTURES LABORATORY}}\\[4pt]
    {\large \textbf{Program: Linked Queue for Student Records}}\\[8pt]
    \rule{\textwidth}{1pt}
\end{center}

\vspace{0.3em}

\section*{Problem Statement}

\noindent
\fcolorbox{codeframe}{codebg}{%
\begin{minipage}{0.97\textwidth}
\vspace{0.4em}
\textbf{Question:}\\[0.3em]
Write a menu-driven program to implement a linked queue for storing student records containing name, roll number, and marks. Implement the following operations:
\begin{enumerate}
    \item \textbf{Enqueue}: Add a student record at the rear of the queue.
    \item \textbf{Dequeue}: Remove a student record from the front of the queue.
    \item \textbf{Display}: Show all student records in the queue.
\end{enumerate}
\vspace{0.2em}
\end{minipage}%
}

\vspace{0.8em}

\section*{Source Code (\texttt{linked-queue.c})}
\lstinputlisting[style=cstyle]{linked-queue.c}

\vspace{0.8em}

\section*{Sample Output}
\begin{lstlisting}[style=terminalstyle]
1. Enqueue
2. Dequeue
3. Display Queue
4. Exit
Choice: 1
Enter name: John
Enter roll no: 101
Enter marks: 85
[+] Enqueued successfully!

1. Enqueue
2. Dequeue
3. Display Queue
4. Exit
Choice: 1
Enter name: Alice
Enter roll no: 102
Enter marks: 92
[+] Enqueued successfully!

1. Enqueue
2. Dequeue
3. Display Queue
4. Exit
Choice: 3

--- Student Queue ---
Roll No: 101 | Name: John | Marks: 85
Roll No: 102 | Name: Alice | Marks: 92

1. Enqueue
2. Dequeue
3. Display Queue
4. Exit
Choice: 2
[+] Dequeued: Roll No: 101 | Name: John | Marks: 85

1. Enqueue
2. Dequeue
3. Display Queue
4. Exit
Choice: 3

--- Student Queue ---
Roll No: 102 | Name: Alice | Marks: 92

1. Enqueue
2. Dequeue
3. Display Queue
4. Exit
Choice: 2
[+] Dequeued: Roll No: 102 | Name: Alice | Marks: 92

1. Enqueue
2. Dequeue
3. Display Queue
4. Exit
Choice: 2
[-] Queue Underflow! No records to dequeue.

1. Enqueue
2. Dequeue
3. Display Queue
4. Exit
Choice: 4
Exiting program...
\end{lstlisting}

\end{document}
